% FORMALIZACION_NUEVA of the lift of N39, compatible with (rho9,q9),
% T_d and Upd_t; preexisting remainder/quotient and orientation architecture.
% Separate research fragment. Does not select the CKM N42 rules.
\subsubsection{Carry transport and integer holonomy of the arithmetic crossing}
\label{subsec:transporte-entero-cruce}

The modular reading of the crossing admits an integral realization retaining
the information accumulated while crossing the edge. This realization
is developed here through transport quotients. Its domain is a lifted
chart of the APP support; the subsequent assignment of a region to an angular
sector retains its own incidence map.

\paragraph{Nonadic chart and lifted transport.}
We use the exact decomposition
\[
 \rho_9(n)=1+((n-1)\bmod9),\qquad
 q_9(n)=\left\lfloor\frac{n-1}{9}\right\rfloor,
 \qquad n=\rho_9(n)+9q_9(n).
\]
A lifted position is
$z=(x,y,q_x,q_y)\in D_9^2\times\mathbb Z^2$, where
$X=x+9q_x$ and $Y=y+9q_y$. The integer increment $u$ of the first
coordinate acts by
\[
 T^x_u(x,y,q_x,q_y)=
 \bigl(\rho_9(x+u),y,q_x+c_9(x;u),q_y\bigr),\qquad
 c_9(x;u)=\left\lfloor\frac{x+u-1}{9}\right\rfloor.
\]
The second coordinate is transported in the same way. Cardinal
advances correspond to $u=\pm1$; the zero increment preserves the state.
The cocycle satisfies
\[
 c_9(x;u+v)=c_9(x;u)+c_9(\rho_9(x+u);v).
\]
Indeed, both sides are the quotient of the unique decomposition of
$x+u+v$ with remainder in $D_9$. Consequently,
$T^x_vT^x_u=T^x_{u+v}$ and $T^x_{-u}=(T^x_u)^{-1}$; the transports
of the two coordinates commute. The state preserves the number of edge
crossings even when the visible position returns.

\paragraph{Lifted sheets and update of records.}
The integral realization is fixed by prolonging the APP arithmetic
operations to the transported coordinates:
\begin{align*}
 \widetilde S(z)&=X+Y=x+y+9(q_x+q_y),\\
 \widetilde P(z)&=XY
 =xy+9(q_xy+q_yx)+81q_xq_y.
\end{align*}
At each position the record of both sheets is preserved:
\[
 \bigl(\rho_9(\widetilde S),q_9(\widetilde S),
       \rho_9(\widetilde P),q_9(\widetilde P)\bigr).
\]
If $e:z\to z'$ is a link and $T$ the selected sheet, its full increment
is obtained from the record before and after transport:
\begin{equation}
 \delta_e\widetilde T=
 \rho_9(\widetilde T(z'))-\rho_9(\widetilde T(z))
 +9\bigl(q_9(\widetilde T(z'))-q_9(\widetilde T(z))\bigr).
 \label{eq:cruce-incremento-registro}
\end{equation}
This identity makes preservation explicit: the residue jump and
quotient increment together reconstruct the variation of the sheet.
Sheet selection, cardinal transport, and record
update are distinct operations. In this chart, the ordered sheet
selection $(T_x,T_y)$ is first composed with $T_d$ and the update
of its quotients; the formula does not postulate that the entire TPK is reduced
to this local representation.

\begin{proposition}[Integral curvature and boundary transport]
\label{prop:cruce-stokes-integral}
Define
$\widetilde A_x=\delta_x\widetilde T_x$ and
$\widetilde A_y=\delta_y\widetilde T_y$ through
\eqref{eq:cruce-incremento-registro}, and let
\[
 \widetilde F_{xy}(X,Y)=
 \widetilde A_x(X,Y)+\widetilde A_y(X+1,Y)
 -\widetilde A_x(X,Y+1)-\widetilde A_y(X,Y).
\]
Homogeneous choices have zero curvature. The mixed choices
$(\widetilde S,\widetilde P)$ and
$(\widetilde P,\widetilde S)$ have curvature
$+1$ and $-1$, respectively, as integers. For the positive boundary of a rectangle
$R_{a,b}$ with integer sides $a,b>0$,
\[
 \widetilde W_{SP}(\partial R_{a,b})=ab,\qquad
 \widetilde W_{PS}(\partial R_{a,b})=-ab.
\]
The result is independent of the base point and the number of crossings
of nonadic seams.
\end{proposition}

\begin{proof}
Exact reconstruction of the two sheets gives
\[
 \delta_x\widetilde S=\delta_y\widetilde S=1,
 \qquad\delta_x\widetilde P=Y,
 \qquad\delta_y\widetilde P=X.
\]
Thus the $SP$ connection is $(1,X)$, with curvature
$1+(X+1)-1-X=1$; the $PS$ connection is $(Y,1)$, with curvature
$Y+1-(Y+1)-1=-1$. For a single sheet, mixed
differences cancel. Summing the curvatures of the $ab$ plaquettes cancels each
interior link with its inverse and preserves the boundary links.
This yields $\pm ab$ in $\mathbb Z$. The transported quotients
ensure the same identity holds across any seam.
\end{proof}

The boundary accumulator can in turn be retained in nonadic form:
$W=r_W+9q_W$. For an integer contribution $w_e$, it is updated by
\[
 r_W'=\rho_9(r_W+w_e),\qquad
 q_W'=q_W+\left\lfloor\frac{r_W+w_e-1}{9}\right\rfloor.
\]
Composition by links provides the same sum as direct integer
evaluation. In this chart, integer zero has coordinates
$(r_W,q_W)=(9,-1)$; this is an arithmetic representation of the accumulator,
without further identification with the distinct oriented zeros of TRIT.

\paragraph{Three operations on orientation.}
Let $\gamma$ be a closed cardinal route. Its time reversal
$\gamma^{-1}$ reverses the order of links and the orientation of each;
by additivity,
\[
 \widetilde W_{SP}(\gamma^{-1})=-\widetilde W_{SP}(\gamma).
\]
Exchange of sheets likewise satisfies
\[
 \widetilde W_{PS}(\gamma)=-\widetilde W_{SP}(\gamma).
\]
To check this on any route, note that the sum of both
connections is the exact increment of
$\widetilde S+\widetilde P=X+Y+XY$; its integral over a loop is zero.

The simultaneous half-turn of both directions, on the other hand, acts
through $J_c(X,Y)=(2c_x-X,2c_y-Y)$, with $2c\in\mathbb Z^2$,
so that it preserves the lattice chart. For a loop,
\[
 \widetilde W_{SP}(J_c\gamma)=\widetilde W_{SP}(\gamma).
\]
Indeed, the closed horizontal contribution is zero, and the vertical one
transforms as
$\sum (2c_x-X)(-\delta Y)=\sum X\delta Y$, since
$\sum\delta Y=0$. A spatial half-turn therefore preserves the
area orientation of this chart. Its combination with time reversal
or sheet exchange is computed by composing the corresponding
operations.

The difference matters for the electronic record: simultaneous reversal
of both displacements every $27$ advances produces, in
windows of $9$, the sequence
\[
 \tau_j=(-1)^{\lfloor j/3\rfloor},\qquad 0\leq j<12,
 \qquad (+,+,+,-,-,-,+,+,+,-,-,-).
\]
This sequence records advance orientations. Holonomy additionally requires
the links traversed, selected sheets, and accumulated
quotients. Transport of electronic orientation is thus preserved
without identifying a half-turn of both axes with time reversal
of a loop.

\paragraph{Application to the marked crossing and sector scope.}
For the rectangular crossing with sides $4$ and $7$, this realization evaluates
exactly
\[
 \widetilde W_{SP}=28=1+9\cdot3,\qquad
 \widetilde W_{SP}(\gamma^{-1})=-28=8+9(-4).
\]
The integer preserves the oriented area; the isolated remainder represents only
its nonadic class. For example, from $(7,1)$ the periodic
representatives of the links sum to $-35$, whereas the carry correction
adds $63$ and recovers $28$.

In the marked incidence of the angular reader, $4$ is
$|H_e\cap P_\pi|$ and $7$ is the marked pivot $p_\varphi$.
The rectangular reading supplies an integer realization of the product
$|H_e\cap P_\pi|\,p_\varphi$ once that crossing has been declared. The
map identifying that region with sector $12$ must preserve
the mark and orientation of the sector system in
Definition~\ref{def:ii09-sistema-sectorial}. Sector assignments
$12,23,13$ and their combinations of $A,h,\Delta^\circ$
are an additional specification beyond the rectangular-holonomy law.
The present result makes the area lift, carry, and
their involutions explicit, with these starting data declared.
